\section{Thảo luận, giới hạn và hướng phát triển}

\subsection{Diễn giải các kết quả chính}
Kết quả cho thấy việc đặt engine làm đơn vị chia dữ liệu và đánh giá giúp hạn chế việc
đếm các window có phụ thuộc thời gian như những quan sát độc lập. Trong
protocol đánh giá, HistGradientBoosting đạt MAE trung bình theo engine thấp nhất
trên đánh giá chéo và được huấn luyện lại trước khi tạo một dự đoán cuối cho mỗi
engine test. Metric cuối vì vậy phản ánh sai số ở điểm cắt của 100 engine test,
không phải sai số của toàn bộ chuỗi theo thời gian.

Phân tích đầu/cuối vòng đời cho thấy nhiều sensor có thay đổi có hệ thống trong
hai đoạn 30 cycles. Tuy nhiên, độ lớn chênh lệch nằm ở đơn vị gốc của từng
sensor, nên không thể dùng các giá trị này để xếp hạng trực tiếp giữa các
sensor. Wilcoxon signed-rank được diễn giải dưới giả định phân phối chênh lệch
đối xứng quanh 0; không nên rút gọn thành kiểm định median cho mọi phân phối.
Với sensor 6, 53 chênh lệch bằng 0 bị loại khỏi phép xếp hạng theo
\texttt{zero\_method=wilcox}, còn 47 chênh lệch dương tham gia kiểm định.
P-value nhỏ không chứng minh quan hệ nhân quả hay giá trị dự báo ngoài mẫu.

PCA--ARIMA cung cấp một cách kiểm tra độc lập về diễn biến của Health Index.
Tỷ lệ bao phủ của khoảng dự đoán 95\% thay đổi đáng kể giữa các engine, từ 21,43\%
đến 97,03\%. ADF trên các đoạn huấn luyện không bác bỏ giả thuyết gốc về nghiệm
đơn vị; do đó, các chuỗi này chưa cung cấp bằng chứng đủ mạnh về tính dừng theo
kiểm định ADF. Sau khi điều chỉnh bậc tự do theo bậc ARIMA, Ljung--Box vẫn cho
thấy p-value dưới 0,05 ở engine 69 và 70. Dự báo Health Index vì vậy chỉ được
giữ như một phân tích phụ, không được dùng để thay thế mô hình RUL.

So sánh với baseline trên cùng prefix và suffix cho thấy ARIMA có MAE trung bình
0,4110, thấp hơn \texttt{drift} (0,7469) và \texttt{last\_value} (0,9478); RMSE tương ứng là
0,5027, 0,9128 và 1,1553. ARIMA chỉ thấp hơn \texttt{drift} ở 7/10 engine và
\texttt{last\_value} ở 8/10 engine, nên kết quả này phù hợp với một lợi thế mô tả trong
subset được chọn hơn là bằng chứng về ưu thế tổng quát. Baseline không có
khoảng dự đoán, vì vậy không được so tỷ lệ bao phủ với ARIMA.

Cảnh báo sớm minh họa trên OOF cho thấy ngưỡng 30 cycles có thể chuyển
dự đoán thành một sự kiện theo engine thay vì đếm từng cycle như các quan sát
độc lập. Tuy nhiên, 49/100 cảnh báo đầu tiên đến sau mốc RUL thật bằng 30;
lead time thấp nhất là 19 cycles. Không có engine hoàn toàn thiếu cảnh báo
không đồng nghĩa mọi cảnh báo đều đủ sớm. Các nhóm trước/đúng/sau ngưỡng chỉ
mô tả thời điểm, chưa phản ánh chi phí hoặc giá trị sử dụng cho bảo trì. Không có engine
âm tính theo định nghĩa sự kiện toàn vòng đời, nên tỷ lệ cảnh báo sai cấp engine
không xác định. Vẫn có thể định nghĩa sai cảnh báo theo từng vùng thời gian,
nhưng không được thay thế âm thầm đại lượng cần ước lượng hiện tại bằng phép đo đó.

Kết quả warning còn chịu ảnh hưởng của model selection: cùng tập OOF đã dùng
để chọn HistGradientBoosting rồi mô tả event. Dự đoán mỗi engine không dùng
engine đó trong fit, nhưng toàn bộ phân tích chưa độc lập với lựa chọn model.
Muốn ước lượng hiệu quả của cả quy trình lựa chọn cần một đánh giá vòng ngoài
riêng; 100\% từng phát cảnh báo không thay thế được đánh giá này.

\subsection{Giới hạn phương pháp}
Thứ nhất, FD001 chỉ gồm một điều kiện vận hành và một dạng suy giảm mô phỏng.
Ma trận bổ sung với window 10 cycles cho thấy MAE transfer từ FD001 tăng lên
60,3215 cycles trên FD002, 47,3949 trên FD003 và 58,5731 trên FD004, trong khi
evaluation nội bộ của các subset nằm khoảng 21,9258--28,3255 cycles. Kết quả này
cho thấy khác biệt phân bố giữa các subset có ảnh hưởng đáng kể trong protocol đang xét; nó không
cho phép suy ra một metric chung cho mọi điều kiện. Thứ hai, test evaluation chỉ có một
điểm cắt cho mỗi engine; cách đánh giá này không mô tả sai số khi mô hình được
gọi lặp lại ở nhiều cycle. Thứ ba, runtime benchmark hiện chỉ đo Pandas trên
dữ liệu nhân bản trong bộ nhớ, với hai lần đo mỗi scale. PySpark mới được
kiểm tra đối chiếu: kết quả tương đương trên input đã xét không phải bằng chứng
về tốc độ hoặc khả năng mở rộng của hệ thống phân tán.
Cuối cùng, dữ liệu không cung cấp chi phí bảo trì, downtime hoặc quy đổi cycle
sang thời gian lịch; vì vậy báo cáo chưa thể chuyển metric RUL thành quyết định
bảo trì.

\subsection{Hướng nghiên cứu tiếp theo}
Khả năng khái quát có thể được khảo sát thêm bằng cách thay đổi độ dài window,
chia dữ liệu theo thời gian trên nhiều engine và dùng đánh giá vòng ngoài cho
quy trình chọn model. Với Health Index, cần huấn luyện transformer ngoài các engine
được dự báo và báo cáo đồng thời tỷ lệ phát hiện lỗi, tỷ lệ bao phủ và khoảng thời gian dự báo.

Kết quả FD002--FD004 hiện đã cung cấp một phép kiểm tra chuyển sang subset khác.
Các bước tiếp theo có thể đánh giá huấn luyện lại trên từng subset, chuẩn hóa theo
operating condition hoặc điều chỉnh theo miền dữ liệu. Mô hình LSTM/Transformer cần một
giả thuyết về biểu diễn thời gian và so sánh cùng protocol với baseline hiện
tại. Khi thiếu cost, downtime và intervention, chưa thể chuyển kết quả RUL
thành chính sách bảo trì.
